\selectlanguage{french}
\chapter*{Résumé}
\addcontentsline{toc}{chapter}{Résumé}

Les véhicules modernes, en particulier les véhicules électriques, reposent sur le bus \gls{can} (Controller Area Network) pour permettre à leurs calculateurs (\gls{ecu}) d'échanger des données. Le protocole \gls{can} a été conçu pour la fiabilité sur un bus fermé, et non pour la sécurité : il ne dispose d'aucune authentification ni chiffrement natif. Dans un système de gestion de batterie (BMS), où le trafic \gls{can} reflète des valeurs critiques pour la sécurité telles que la tension des cellules, l'état de charge et la température, les ingénieurs ont besoin d'une visibilité sur ce trafic qui va au-delà d'un simple enregistrement passif des trames.

Ce rapport présente \textbf{CANvisionNative}, une plateforme intelligente de diagnostic et d'aide à la validation \gls{can} pour les systèmes de gestion de batterie (BMS). Son objectif principal est le diagnostic : offrir à un ingénieur BMS une visibilité continue et unifiée sur le trafic \gls{can}, là où les outils actuels restent fragmentés et statiques. Autour de cet objectif, la plateforme combine surveillance en temps réel, rejeu de journaux et détection d'anomalies fondée sur l'apprentissage automatique au sein d'un même poste de travail. Le système combine un client de bureau Windows, développé en \gls{wpf} selon le patron \gls{mvvm}, avec un backend Python exposant une \gls{api} \gls{rest} construite avec FastAPI. Le backend décode les trames \gls{can} brutes à l'aide de fichiers \gls{dbc} spécifiques à chaque véhicule, extrait des caractéristiques temporelles et de contenu, puis évalue chaque trame par la fusion de quatre couches d'analyse : une couche temporelle, une couche de continuité de charge utile, une couche d'apprentissage automatique (Isolation Forest) propre à chaque véhicule, et une couche de persistance temporelle. Les anomalies détectées sont expliquées en langage naturel grâce à un module d'explication combinant un raisonnement à base de règles et une chaîne de modèles de langage (\gls{llm}).

La plateforme prend en charge trois modes de fonctionnement : le rejeu de journaux enregistrés (\gls{csv}, \gls{mf4}, candump, Vector \gls{can}), la simulation en direct, et la capture matérielle en direct via un adaptateur ELM327 ou SocketCAN. Sa couche d'apprentissage automatique a été entraînée sur plus de sept millions de trames \gls{can} provenant de seize sources de véhicules et de jeux de données, et validée sur un enregistrement réel d'une heure d'une Kia EV6 contenant 97~741 trames.

Le projet a suivi une méthodologie d'ingénierie fondée sur la preuve : chaque affirmation sur le comportement du système a été vérifiée par des mesures reproductibles avant d'être retenue. Cette rigueur a permis d'aboutir à un pipeline de détection entièrement corrigé, dont la version finale réduit le volume d'alertes de 25~\% et élimine totalement la saturation du score d'apprentissage automatique, sans assouplir le seuil de détection.

Ce rapport documente les besoins, l'architecture, l'implémentation et la validation de CANvisionNative, et montre qu'une plateforme de diagnostic et d'aide à la validation \gls{can}, multicouche, explicable et de classe bureautique, peut être construite et rigoureusement vérifiée pour des données réelles de bus \gls{can} de véhicule, en utilisant du matériel courant et des jeux de données ouverts.

\vspace{0.5cm}
\noindent\textbf{Mots-clés :} Bus \gls{can}, Système de gestion de batterie, Diagnostic \gls{can}, Véhicule électrique, Apprentissage automatique, Isolation Forest, Détection d'anomalies, Modèle de langage, Ingénierie de validation automobile.

\selectlanguage{english}
\cleardoublepage
